% GENERATED FILE. Edit canonical lessons, then run npm run generate:books.
% canonical-source-hash: fnv1a64:9ac80128f494b48a
% canonical-lessons: SA-C554-sambhavati, SA-C554-dirghadarsi, SA-C554-abhipraiti, SA-C554-ayatih, SA-C554-arambhah

\chapter{Looking Ahead}
\label{ch:sa-looking-ahead}
\hlchaptermodality{554}

\begin{chapteropening}
\textbf{By the end of this chapter:} \emph{I can say what I intend and what the future may bring.}

Complete the last lesson of chapter 554: I can say what I intend and what the future may bring.
\end{chapteropening}

\section[\texorpdfstring{\sk{सम्भवति} (sambhavati)}{सम्भवति (sambhavati)}]{\sk{सम्भवति} (sambhavati) — is possible; comes about}

\label{lesson:SA-C554-sambhavati}



\begin{quote}
Before the new one: say the Sanskrit for a reversal, then the Sanskrit for a reproach.
\end{quote}

\subsection*{You'll want to know: \sk{सम्भवति}}
\textbf{\sk{सम्भवति}} — \emph{sambhavati} — \textquotedblleft{}is possible; comes about\textquotedblright{}.

\begin{quote}
\sk{यदि} \sk{श्वः} \sk{वर्षा} \sk{सम्भवति}, \sk{तर्हि} \sk{छत्रं} \sk{नयतु।}
\end{quote}

\begin{quote}
\emph{yadi śvaḥ varṣā sambhavati, tarhi chatraṃ nayatu.}
\end{quote}

\begin{quote}
If rain is possible tomorrow, take an umbrella.
\end{quote}

Past (it …): \sk{उत्सवः} \sk{सम्भूतवान्}, \emph{utsavaḥ sambhūtavān}.

\begin{cousinweb}
\textbf{\sk{सम्भवति}} (\emph{sambhavati}) is built on \emph{sam}, \textquotedblleft{}together\textquotedblright{}, and the root √bhū, \textquotedblleft{}to be, to become\textquotedblright{}, which also gives \emph{bhavati}; MW gives \emph{sambhū} as to come together, to arise, to happen, and to be possible; its past participle is \emph{sambhūta} (MW).
\end{cousinweb}

\begin{grammarlens}[title={What you've built}]
One more word for saying what happens if.
\end{grammarlens}

\subsection*{Your turn}
\begin{itemize}
\raggedright
  \item \emph{Say it:} \emph{sambhavati}
  \item \emph{Say it:} \emph{sambhavati}, once more
  \item \emph{Say it:} \emph{yadi śvaḥ varṣā sambhavati, tarhi chatraṃ nayatu.} — If rain is possible tomorrow, take an umbrella.
  \item \emph{Recall:} say the Sanskrit for a reversal, then the Sanskrit for a reproach, then say \emph{sambhavati} again
\end{itemize}

\begin{tcolorbox}[breakable,colback=teal!4,colframe=teal!35!black,title={Before you move on}]
What does \sk{सम्भवति} mean? (\textquotedblleft{}Is possible; comes about\textquotedblright{}.) Say it once more.
\end{tcolorbox}
\section[\texorpdfstring{\sk{दीर्घदर्शी} (dīrghadarśī)}{दीर्घदर्शी (dīrghadarśī)}]{\sk{दीर्घदर्शी} (dīrghadarśī) — far-sighted}

\label{lesson:SA-C554-dirghadarsi}



\begin{quote}
Before the new one: say the Sanskrit for a reproach, then the Sanskrit for is possible.
\end{quote}

\subsection*{You'll want to know: \sk{दीर्घदर्शी}}
\textbf{\sk{दीर्घदर्शी}} — \emph{dīrghadarśī} — \textquotedblleft{}far-sighted\textquotedblright{}.

\begin{quote}
\sk{यदि} \sk{नायकः} \sk{दीर्घदर्शी} \sk{अस्ति}, \sk{तर्हि} \sk{देशः} \sk{सुरक्षितः} \sk{भवति।}
\end{quote}

\begin{quote}
\emph{yadi nāyakaḥ dīrghadarśī asti, tarhi deśaḥ surakṣitaḥ bhavati.}
\end{quote}

\begin{quote}
If the leader is far-sighted, the country is safe.
\end{quote}

\begin{cousinweb}
\textbf{\sk{दीर्घदर्शी}} (\emph{dīrghadarśī}) is built from \emph{dīrgha}, \textquotedblleft{}long, far\textquotedblright{}, and \emph{darśin}, \textquotedblleft{}seeing\textquotedblright{}, from the root √dṛś, \textquotedblleft{}to see\textquotedblright{}; MW gives \emph{dīrghadarśin} as the same as \emph{dīrghadarśa}, far-sighted (MW).
\end{cousinweb}

\begin{grammarlens}[title={What you've built}]
One more word for saying what happens if.
\end{grammarlens}

\subsection*{Your turn}
\begin{itemize}
\raggedright
  \item \emph{Say it:} \emph{dīrghadarśī}
  \item \emph{Say it:} \emph{dīrghadarśī}, once more
  \item \emph{Say it:} \emph{yadi nāyakaḥ dīrghadarśī asti, tarhi deśaḥ surakṣitaḥ bhavati.} — If the leader is far-sighted, the country is safe.
  \item \emph{Recall:} say the Sanskrit for a reproach, then the Sanskrit for is possible, then say \emph{dīrghadarśī} again
\end{itemize}

\begin{tcolorbox}[breakable,colback=teal!4,colframe=teal!35!black,title={Before you move on}]
What does \sk{दीर्घदर्शी} mean? (\textquotedblleft{}Far-sighted\textquotedblright{}.) Say it once more.
\end{tcolorbox}
\section[\texorpdfstring{\sk{अभिप्रैति} (abhipraiti)}{अभिप्रैति (abhipraiti)}]{\sk{अभिप्रैति} (abhipraiti) — intends}

\label{lesson:SA-C554-abhipraiti}



\begin{quote}
Before the new one: say the Sanskrit for is possible, then the Sanskrit for far-sighted.
\end{quote}

\subsection*{You'll want to know: \sk{अभिप्रैति}}
\textbf{\sk{अभिप्रैति}} — \emph{abhipraiti} — \textquotedblleft{}intends\textquotedblright{}.

\begin{quote}
\sk{यदि} \sk{भवान्} \sk{गन्तुम्} \sk{अभिप्रैति}, \sk{तर्हि} \sk{अहम्} \sk{अपि} \sk{आगच्छामि।}
\end{quote}

\begin{quote}
\emph{yadi bhavān gantum abhipraiti, tarhi aham api āgacchāmi.}
\end{quote}

\begin{quote}
If you intend to go, I will come too.
\end{quote}

Past (he …): \sk{सः} \sk{गन्तुम्} \sk{अभिप्रेतवान्}, \emph{saḥ gantum abhipretavān}.

\begin{cousinweb}
\textbf{\sk{अभिप्रैति}} (\emph{abhipraiti}) is built on \emph{abhi}, \textquotedblleft{}towards\textquotedblright{}, \emph{pra}, \textquotedblleft{}forth\textquotedblright{}, and the root √i, \textquotedblleft{}to go\textquotedblright{}; MW gives \emph{abhipre} as to approach in mind, to aim at, to intend; its past participle is \emph{abhipreta} (MW).
\end{cousinweb}

\begin{grammarlens}[title={What you've built}]
One more word for saying what happens if.
\end{grammarlens}

\subsection*{Your turn}
\begin{itemize}
\raggedright
  \item \emph{Say it:} \emph{abhipraiti}
  \item \emph{Say it:} \emph{abhipraiti}, once more
  \item \emph{Say it:} \emph{yadi bhavān gantum abhipraiti, tarhi aham api āgacchāmi.} — If you intend to go, I will come too.
  \item \emph{Recall:} say the Sanskrit for is possible, then the Sanskrit for far-sighted, then say \emph{abhipraiti} again
\end{itemize}

\begin{tcolorbox}[breakable,colback=teal!4,colframe=teal!35!black,title={Before you move on}]
What does \sk{अभिप्रैति} mean? (\textquotedblleft{}Intends\textquotedblright{}.) Say it once more.
\end{tcolorbox}
\section[\texorpdfstring{\sk{आयतिः} (āyatiḥ)}{आयतिः (āyatiḥ)}]{\sk{आयतिः} (āyatiḥ) — the time ahead}

\label{lesson:SA-C554-ayatih}



\begin{quote}
Before the new one: say the Sanskrit for far-sighted, then the Sanskrit for intends.
\end{quote}

\subsection*{You'll want to know: \sk{आयतिः}}
\textbf{\sk{आयतिः}} — \emph{āyatiḥ} — \textquotedblleft{}the time ahead\textquotedblright{}.

\begin{quote}
\sk{यदि} \sk{आयतिः} \sk{शुभा} \sk{अस्ति}, \sk{तर्हि} \sk{वयं} \sk{प्रसन्नाः} \sk{स्मः।}
\end{quote}

\begin{quote}
\emph{yadi āyatiḥ śubhā asti, tarhi vayaṃ prasannāḥ smaḥ.}
\end{quote}

\begin{quote}
If the future looks bright, we are glad.
\end{quote}

A feminine noun.

\begin{cousinweb}
\textbf{\sk{आयतिः}} (\emph{āyatiḥ}) is built on \emph{ā} and the root √yam, \textquotedblleft{}to hold, to stretch\textquotedblright{}: a stretching out; MW gives \emph{āyati} as extension, and following or future time (MW).
\end{cousinweb}

\begin{grammarlens}[title={What you've built}]
One more word for saying what happens if.
\end{grammarlens}

\subsection*{Your turn}
\begin{itemize}
\raggedright
  \item \emph{Say it:} \emph{āyatiḥ}
  \item \emph{Say it:} \emph{āyatiḥ}, once more
  \item \emph{Say it:} \emph{yadi āyatiḥ śubhā asti, tarhi vayaṃ prasannāḥ smaḥ.} — If the future looks bright, we are glad.
  \item \emph{Recall:} say the Sanskrit for far-sighted, then the Sanskrit for intends, then say \emph{āyatiḥ} again
\end{itemize}

\begin{tcolorbox}[breakable,colback=teal!4,colframe=teal!35!black,title={Before you move on}]
What does \sk{आयतिः} mean? (\textquotedblleft{}The time ahead\textquotedblright{}.) Say it once more.
\end{tcolorbox}
\section[\texorpdfstring{\sk{आरम्भः} (ārambhaḥ)}{आरम्भः (ārambhaḥ)}]{\sk{आरम्भः} (ārambhaḥ) — a start}

\label{lesson:SA-C554-arambhah}



\begin{quote}
Before the new one: say the Sanskrit for intends, then the Sanskrit for the time ahead.
\end{quote}

\subsection*{You'll want to know: \sk{आरम्भः}}
\textbf{\sk{आरम्भः}} — \emph{ārambhaḥ} — \textquotedblleft{}a start\textquotedblright{}.

\begin{quote}
\sk{यदि} \sk{आरम्भः} \sk{उत्तमः} \sk{अस्ति}, \sk{तर्हि} \sk{फलम्} \sk{अपि} \sk{उत्तमं} \sk{भविष्यति।}
\end{quote}

\begin{quote}
\emph{yadi ārambhaḥ uttamaḥ asti, tarhi phalam api uttamaṃ bhaviṣyati.}
\end{quote}

\begin{quote}
If the start is good, the result will be good too.
\end{quote}

A masculine noun.

\begin{cousinweb}
\textbf{\sk{आरम्भः}} (\emph{ārambhaḥ}) is built on \emph{ā} and the root √rabh, \textquotedblleft{}to take hold of\textquotedblright{}, which also gives \emph{ārabhate}, \textquotedblleft{}begins\textquotedblright{}; MW gives \emph{ārambha} as undertaking, beginning (MW).
\end{cousinweb}

\begin{grammarlens}[title={What you've built}]
One more word for saying what happens if.
\end{grammarlens}

\subsection*{Your turn}
\begin{itemize}
\raggedright
  \item \emph{Say it:} \emph{ārambhaḥ}
  \item \emph{Say it:} \emph{ārambhaḥ}, once more
  \item \emph{Say it:} \emph{yadi ārambhaḥ uttamaḥ asti, tarhi phalam api uttamaṃ bhaviṣyati.} — If the start is good, the result will be good too.
  \item \emph{Recall:} say the Sanskrit for intends, then the Sanskrit for the time ahead, then say \emph{ārambhaḥ} again
\end{itemize}

\begin{tcolorbox}[breakable,colback=teal!4,colframe=teal!35!black,title={Before you move on}]
What does \sk{आरम्भः} mean? (\textquotedblleft{}A start\textquotedblright{}.) Say it once more.
\end{tcolorbox}
